\documentclass[14pt,a4paper]{report}
\usepackage[utf8]{inputenc}
\usepackage[french]{babel}
\usepackage{graphicx}
\usepackage{hyperref}
\usepackage{listings}
\usepackage{color}
\usepackage{enumitem}
\usepackage{fancyhdr}
\usepackage{geometry}
\usepackage{booktabs}
\usepackage{multirow}
\usepackage{array}
\usepackage{siunitx}

% Configuration de la géométrie de la page
\geometry{
    a4paper,
    top=2.5cm,
    bottom=2.5cm,
    left=2.5cm,
    right=2.5cm
}

% Configuration des en-têtes et pieds de page
\pagestyle{fancy}
\fancyhf{}
\rhead{Mercato}
\lhead{Étude de Faisabilité}
\rfoot{Page \thepage}

\begin{document}

\begin{titlepage}
    \begin{center}
        {\large REPUBLIQUE DU CAMEROUN \hfill REPUBLIC OF CAMEROON}\\
        {\small Paix-Travail-Patrie \hfill Peace-Work-Fatherland}\\[0.5cm]
        {\large UNIVERSITE DE YAOUNDE 1 \hfill UNIVERSITY OF YAOUNDE 1}\\
        {\large FACULTE DES SCIENCES \hfill FACULTY OF SCIENCE}\\
        {\large Département d'Informatique \hfill Department of computer science}\\
        {\small B.P.812 Yaoundé \hfill PO.BOX 812 Yaoundé}\\
        {\small Tél /Fax : (+237) 22 23 44 96 \hfill Tel /Fax : (+237) 22 23 44 96}\\[2cm]

        \begin{center}
            {\large \textbf{Filière} : Informatique \hspace{2cm} \textbf{Niveau} : M1}\\
            {\large \textbf{Unité d'enseignement (UE)} : INF 4027 : Génie Logiciel}\\[2cm]

            \framebox{
                \begin{minipage}{0.8\textwidth}
                    \begin{center}
                        \huge \textbf{MERCATO}\\[0.5cm]
                        \Large Étude de Faisabilité
                    \end{center}
                \end{minipage}
            }\\[2cm]

            \textbf{Rédigé par} :\\[0.5cm]
            \begin{tabular}{|l|l|}
                \hline
                \textbf{Noms \& Prénoms} & \textbf{Matricule} \\
                \hline
                MAHAMAT SALEH MAHAMAT & 21T2423 \\
                \hline
                SOKOUDJOU CHENDOU CHRISTIAN MANUEL & 21T2396 \\
                \hline
                ZOGO ABOUMA ZOZINE ACHAIRE & 18N2824 \\
                \hline
                FEUJO TSAGMO SIMPLICE JORDA & 18T2897 \\
                \hline
            \end{tabular}\\[2cm]

            \textbf{Supervisé par} :\\
            Pr. Roger ATSA ETOUNDI\\[1cm]

            \textbf{Année universitaire} : 2024 - 2025
        \end{center}
    \end{center}
\end{titlepage}

\tableofcontents
\newpage

\chapter{Analyse du Contexte}

\section{Description du Projet}
Mercato est une application de gestion de caisse destinée à un usage interne, avec les caractéristiques suivantes :
\begin{itemize}
    \item Utilisation exclusive par les caissières et l'administrateur
    \item Déploiement en local sur serveur dédié
    \item Gestion des ventes et des stocks
    \item Reporting et suivi des transactions
\end{itemize}

\section{Utilisateurs Cibles}
\begin{itemize}
    \item Caissières (utilisateurs principaux)
    \item Administrateur système
    \item Responsable du magasin
\end{itemize}

\section{Analyse des Besoins Utilisateurs}

\subsection{Besoins des Caissières}
\begin{itemize}
    \item \textbf{Gestion des Ventes}
        \begin{itemize}
            \item Interface simple et intuitive pour la saisie des ventes
            \item Scan des codes-barres des produits
            \item Calcul automatique du total et de la monnaie à rendre
            \item Impression des tickets de caisse
            \item Annulation/modification des transactions en cours
        \end{itemize}
    \item \textbf{Gestion des Paiements}
        \begin{itemize}
            \item Paiement en espèces
            \item Suivi du fond de caisse
            \item Clôture de caisse en fin de journée
        \end{itemize}
    \item \textbf{Support Opérationnel}
        \begin{itemize}
            \item Vérification rapide des prix
            \item Consultation des stocks disponibles
            \item Système d'aide contextuelle
            \item Gestion des retours et remboursements
        \end{itemize}
\end{itemize}

\subsection{Besoins de l'Administrateur}
\begin{itemize}
    \item \textbf{Gestion des Utilisateurs}
        \begin{itemize}
            \item Création/modification des comptes caissières
            \item Attribution des droits d'accès
            \item Réinitialisation des mots de passe
        \end{itemize}
    \item \textbf{Configuration Système}
        \begin{itemize}
            \item Paramétrage des caisses
            \item Gestion des imprimantes
            \item Configuration des backups
            \item Maintenance du système
        \end{itemize}
    \item \textbf{Gestion des Produits}
        \begin{itemize}
            \item Ajout/modification des produits
            \item Gestion des prix et promotions
            \item Mise à jour des stocks
        \end{itemize}
\end{itemize}

\subsection{Besoins du Responsable}
\begin{itemize}
    \item \textbf{Reporting}
        \begin{itemize}
            \item Rapport des ventes journalières
            \item Suivi du chiffre d'affaires
            \item Analyse des tendances
            \item États des stocks
        \end{itemize}
    \item \textbf{Supervision}
        \begin{itemize}
            \item Suivi des performances des caissières
            \item Contrôle des annulations et remboursements
            \item Validation des opérations exceptionnelles
        \end{itemize}
    \item \textbf{Gestion des Stocks}
        \begin{itemize}
            \item Alerte de rupture de stock
            \item Suivi des mouvements de stock
            \item Inventaire
        \end{itemize}
\end{itemize}

\subsection{Exigences Non Fonctionnelles}
\begin{itemize}
    \item \textbf{Performance}
        \begin{itemize}
            \item Temps de réponse < 2 secondes
            \item Disponibilité système 99.9\%
            \item Sauvegarde quotidienne des données
        \end{itemize}
    \item \textbf{Sécurité}
        \begin{itemize}
            \item Authentification sécurisée
            \item Traçabilité des opérations
            \item Protection des données sensibles
        \end{itemize}
    \item \textbf{Ergonomie}
        \begin{itemize}
            \item Interface adaptée à l'utilisation tactile
            \item Temps de formation < 2 jours
            \item Documentation en français
        \end{itemize}
\end{itemize}

\chapter{Faisabilité Technique}

\section{Architecture Technique}
\subsection{Infrastructure Locale}
\begin{itemize}
    \item Serveur local dédié
    \item Réseau local (LAN)
    \item Postes clients (ordinateurs des caisses)
    \item Système de sauvegarde local
\end{itemize}

\subsection{Stack Technologique}
\begin{itemize}
    \item Frontend : React.js (interface simple et intuitive)
    \item Backend : Laravel 10 (API RESTful)
    \item Base de données : MySQL (plus adapté pour usage local)
    \item Serveur : Ubuntu Server LTS
\end{itemize}

\chapter{Méthodologie d'Estimation des Coûts}

\section{Approche d'Estimation}
\begin{itemize}
    \item Méthode COCOMO II pour l'estimation du développement
    \item Analyse comparative des salaires locaux
    \item Consultation des fournisseurs locaux pour le matériel
    \item Prise en compte du contexte PME
\end{itemize}

\chapter{Analyse Financière Détaillée}

\section{Coûts de Développement}
\begin{table}[h]
    \centering
    \begin{tabular}{|l|r|r|p{4cm}|}
        \hline
        \textbf{Poste} & \textbf{Coût Mensuel (XAF)} & \textbf{Durée} & \textbf{Justification} \\
        \hline
        Développeur Senior & 400.000 & 6 mois & Base sur le salaire moyen local + expérience \\
        \hline
        Développeur Junior & 250.000 & 6 mois & Salaire entrée de marché \\
        \hline
        Designer UI & 200.000 & 3 mois & Mission ponctuelle \\
        \hline
        \textbf{Total RH} & \multicolumn{3}{r|}{\textbf{5.100.000 XAF}} \\
        \hline
    \end{tabular}
    \caption{Coûts des Ressources Humaines (Phase développement)}
\end{table}

\section{Coûts d'Infrastructure}
\begin{table}[h]
    \centering
    \begin{tabular}{|l|r|p{6cm}|}
        \hline
        \textbf{Équipement} & \textbf{Coût (XAF)} & \textbf{Justification} \\
        \hline
        Serveur Dell PowerEdge T150 & 1.500.000 & Serveur entry-level adapté aux PME \\
        \hline
        Switch réseau & 150.000 & Connexion des postes clients \\
        \hline
        UPS & 300.000 & Protection électrique \\
        \hline
        Installation réseau & 200.000 & Câblage et configuration \\
        \hline
        \textbf{Total Infrastructure} & \textbf{2.150.000} & \textit{Investissement unique} \\
        \hline
    \end{tabular}
    \caption{Coûts d'Infrastructure Initiale}
\end{table}

\section{Coûts Opérationnels Mensuels}
\begin{table}[h]
    \centering
    \begin{tabular}{|l|r|p{6cm}|}
        \hline
        \textbf{Poste} & \textbf{Coût Mensuel (XAF)} & \textbf{Justification} \\
        \hline
        Maintenance technique & 100.000 & Support technique à temps partiel \\
        \hline
        Électricité (serveur) & 30.000 & Basé sur la consommation moyenne \\
        \hline
        Sauvegardes & 20.000 & Stockage externe mensuel \\
        \hline
        \textbf{Total Opérationnel} & \textbf{150.000} & \textit{Coût mensuel récurrent} \\
        \hline
    \end{tabular}
    \caption{Coûts Opérationnels}
\end{table}

\section{Formation et Support}
\begin{table}[h]
    \centering
    \begin{tabular}{|l|r|p{6cm}|}
        \hline
        \textbf{Activité} & \textbf{Coût (XAF)} & \textbf{Détails} \\
        \hline
        Formation initiale & 300.000 & Formation des caissières (5 jours) \\
        \hline
        Documentation & 100.000 & Manuels et guides d'utilisation \\
        \hline
        \textbf{Total Formation} & \textbf{400.000} & \textit{Coût unique} \\
        \hline
    \end{tabular}
    \caption{Coûts de Formation}
\end{table}

\chapter{Résumé des Investissements}

\section{Investissement Initial Total}
\begin{itemize}
    \item Développement : 5.100.000 XAF
    \item Infrastructure : 2.150.000 XAF
    \item Formation : 400.000 XAF
    \item \textbf{Total} : 7.650.000 XAF
\end{itemize}

\section{Coûts Mensuels Récurrents}
\begin{itemize}
    \item Maintenance et opérations : 150.000 XAF
    \item \textbf{Total Annuel} : 1.800.000 XAF
\end{itemize}

\end{document}
